\begin{table*}[t]
\centering
\scriptsize
\renewcommand{\arraystretch}{0.88}
\begin{tabularx}{\linewidth}{>{\raggedright\arraybackslash}p{0.18\linewidth}c c c c >{\raggedright\arraybackslash}p{0.15\linewidth}X}
\toprule
\textbf{Threat class} & \textbf{Corpus} & \textbf{Query} & \textbf{Retr.} & \textbf{Gen.} & \textbf{Objective} & \textbf{Boundary} \\
\midrule
Dense-retriever poisoning \citep{poisoningRetrievalCorpora2023} & write & train set & white/gray & no & retrieve bad passages & Not always end-to-end. \\
RAG knowledge corruption \citep{poisonedRAG2024,badRAG2024} & write & target topic & black/gray & black & answer steering & Stealth/dynamics vary. \\
Backdoor-style RAG \citep{trojanRAG2024,phantom2024} & write & trigger & gray & black & latent behavior & Normal queries may pass. \\
Indirect prompt injection \citep{aip2025,rrag005} & write/share & natural & black & black & instruction hijack & Must overcome retrieval. \\
Opinion/fairness manipulation \citep{rrag104,rrag162} & write & group/topic & black/gray & black & bias steering & QA metrics miss social harm. \\
Defense setting \citep{certifiedRAG2024,rrag004,rrag015,rrag128} & bounded & varies & assumed & black/gray & detect/demote/trace & Depends on budget/visibility. \\
\bottomrule
\end{tabularx}
\caption{Attacker-capability matrix for adversarial RAG. The same ``attack success rate'' can mean different things depending on corpus write access, query control, retriever visibility, generator access, and objective.}
\label{tab:attacker-capability}
\end{table*}
